\section{A barrier to deleting both numerical arguments}
\label{sec:finite-menu}

The freedom to choose a witness and the freedom to choose a specialization
target cannot, in general, be removed together.  The result below concerns
the consequences of a specified native calculus.  It is not a lower bound
on the size of every possible iPRA presentation, and it does not exclude a
different interpretation with a different numerical domain.

\subsection{The restricted calculus}

Numerical terms are
\[
 t ::= x\mid 0\mid S(t)\mid I(n,a;u.h),
\]
where the last constructor binds only $u$ in $h$.  Its computation equations
are
\begin{align*}
 I(0,a;u.h)&=a,\\
 I(Sn,a;u.h)&=h[I(n,a;u.h)/u].
\end{align*}
Formulas are equations, $\mathsf Q_y B(y)$, and
$\mathsf D_{x,y}B(x,y)$.  Substitution is capture avoiding, and all variables
displayed as fresh below satisfy the usual eigenvariable conditions.

Fix once and for all a finite family
$F=(f_1,\ldots,f_m)$ of numerical operations, each specified by a fixed
scalar term $f_i(x_1,\ldots,x_{r_i})$.  The arities $r_i$ are fixed and may
be zero.  For a numerical parameter scope $V$, let $\mathcal T_F(V)$ be the
smallest family containing the parameters in $V$ and all standard numerals,
and closed under composition with the operations in $F$.  This is a
restriction on specialization targets; the general numerical term language
still contains arbitrary iterator bodies.  The family $F$ is fixed for the
whole calculus, not chosen anew for each proof.

Let $\mathcal N_F$ retain hypothesis, both computation equations, and the
following rules of the native calculus:
\begin{align*}
 &\frac{\Gamma\vdash s=t\qquad\Gamma\vdash A(s)}
        {\Gamma\vdash A(t)}
       &&\text{(transport in any native formula template)},\\[2mm]
 &\frac{\Gamma\vdash\mathsf Q_y B(y)\qquad\Gamma\vdash B(0)}
        {\Gamma\vdash B(t)}
       &&\bigl(t\in\mathcal T_F(V)\bigr),\\[2mm]
 &\frac{\Gamma\vdash\mathsf D_{x,y}B(x,y)
        \qquad\Gamma,\mathsf Q_yB(b,y)\vdash A}
        {\Gamma\vdash A}
       &&\bigl(b\text{ fresh}\bigr),\\[2mm]
 &\frac{\Gamma,B(w,0)\vdash B(w,Sz)}
        {\Gamma\vdash\mathsf D_{x,y}B(x,y)}
       &&\bigl(z\text{ fresh},\quad w=0\text{ or }w\in V\bigr),\\[2mm]
 &\frac{\Gamma,E(z)\vdash E(Sz)}
        {\Gamma\vdash\mathsf Q_yE(y)}
       &&\bigl(z\text{ fresh},\quad E\text{ a literal equation}\bigr).
\end{align*}
The fourth rule is introduction by numerical cases, and the fifth is
equation progress.  Contexts are lists of retained assumptions, so an
assumption may be unused.  Restricting the fourth rule further to the
newest parameter, with $0$ in an empty scope, gives a subsystem of
$\mathcal N_F$.

No unrestricted structural rule substituting arbitrary numerical terms
for the free parameters of a finished derivation is added to this
restricted calculus, nor is its admissibility assumed. Such a rule would
restore precisely a freedom removed by the restrictions. Capture-avoiding
substitution within each displayed inference still has its usual meaning.

\paragraph{Theorem (finite-menu obstruction).}
\emph{For every fixed finite family $F$, the calculus $\mathcal N_F$ fails
to derive a closed, $\mathsf Q$-free formula that the unrestricted native
calculus derives.  In particular, allowing a fixed finite menu of numerical
operations as specialization targets does not repair the simultaneous
restriction of introduction witnesses to $0$ and parameters.}

The proof gives an explicit form of the missing formula.  A direct
majorization lemma first ensures that all functions used in its construction
are scalar terms.  Thus this argument does not need to assume the general
compilation of every primitive recursive function into the scalar language.

\subsection{A scalar majorant for a finite menu}

In this subsection, inequalities and induction are used in the ambient
theory $\mathrm{PA}$, with definitional extensions by fixed primitive
recursive functions.  They are not additional rules of $\mathcal N_F$.
Write
\[
 n+r:=I(n,r;v.Sv).
\]
PA proves the usual laws of this addition term.

\paragraph{Lemma (structural majorization).}
\emph{For every fixed scalar term $t$, one can construct a unary scalar
term $B_t(n)$ such that PA proves all three assertions:
\begin{enumerate}
 \item $n<B_t(n)$;
 \item $n<n'\ \Longrightarrow\ B_t(n)<B_t(n')$;
 \item if every free numerical argument of $t$ is at most $n$, then
       $t\le B_t(n)$.
\end{enumerate}}

\begin{proof}
Construct the terms by structural recursion on $t$:
\begin{align*}
 B_0(n)&=Sn,& B_{x_j}(n)&=Sn,\\
 B_{St}(n)&=S(B_t(n)),\\
 B_{I(c,a;u.h)}(n)
   &=I\bigl(B_c(n),B_a(n);r.B_h(n+r)\bigr).
\end{align*}
When constructing $B_h$, regard the iterator variable $u$ as one more free
argument of $h$.  The variable $r$ in the last display is fresh.  Each
$B_t$ has only its displayed parameter $n$ free.

The zero, variable, and successor cases are immediate.  For the iterator
case, fix $n$ and put
\[
 N=B_c(n),\qquad A=B_a(n),\qquad
 K_n(r)=B_h(n+r),\qquad W_j=I(j,A;r.K_n(r)).
\]
The inductive properties of $B_h$ give $K_n(r)>r$, so PA proves that
$W_j$ is strictly increasing in $j$ and $W_j\ge A>n$.

Suppose the original free arguments are at most $n$.  Their values in the
original counter and seed satisfy $c\le N$ and $a\le A$.  Compare the
original iteration with $W_j$ up to its actual counter $c$.  If the
original accumulator after $j$ steps is at most $W_j$, then every argument
of $h$, including that accumulator, is at most $n+W_j$.  The inductive
majorization property of $B_h$ therefore bounds the next original
accumulator by
\[
 B_h(n+W_j)=W_{j+1}.
\]
PA induction on $j$ proves that the original result is at most
$W_c\le W_N=B_{I(c,a;u.h)}(n)$.

For strict monotonicity in $n$, take $n<n'$.  Both the seed $B_a(n)$ and
the counter $B_c(n)$ are strictly smaller than their counterparts at $n'$.
The update $K_n(r)=B_h(n+r)$ is strictly increasing in $r$ and increases
when $n$ increases.  PA induction first compares the two iterations for
the smaller number of steps, preserving the strict inequality between
their accumulators.  The additional steps of the iteration at $n'$ can
only increase its value.  This proves strict monotonicity of the displayed
bound, as well as its inflationary property.

This is a metatheoretic induction on a fixed finite syntax tree.  Every
induction on numerical values used within it is a PA induction, so the
resulting inequalities hold in every model of PA, including nonstandard
ones.
\end{proof}

Apply the lemma to the finitely many terms specifying $F$, and define the
unary scalar term
\[
 H(n)=n+1+\sum_{i=1}^{m}B_{f_i}(n).
\]
The sum is a fixed finite term built from the addition macro; for $m=0$
it is $0$.  PA proves that $H$ is strictly increasing, that $H(n)>n$, and
that
\begin{equation}
 \bigwedge_{j=1}^{r_i}x_j\le n
 \quad\Longrightarrow\quad
 f_i(x_1,\ldots,x_{r_i})<H(n).
 \label{eq:finite-menu-majorant}
\end{equation}
Indeed, the majorization lemma bounds the left side's value by
$B_{f_i}(n)$, and the additional summand $n+1$ makes the final inequality
strict.  The same argument covers a nullary operation.

Define two more scalar terms:
\[
 J(k,n)=I(k,n;u.H(u)),\qquad
 g(n)=J(n,n).
\]
Their defining equations give
$J(0,n)=n$ and $J(Sk,n)=H(J(k,n))$.  Since $H(x)>x$, PA proves
\begin{equation}
 r<s\quad\Longrightarrow\quad J(r,n)<J(s,n).
 \label{eq:finite-menu-orbit}
\end{equation}
There is no uniform evaluator for all primitive recursive programs hidden
here: $H$, $J$, and $g$ are three fixed finite scalar terms constructed from
the fixed menu $F$.

\subsection{An external cut closed under every permitted target}

Work in the usual classical metatheory.  By compactness, choose a
nonstandard model $M\models\mathrm{PA}$ and an element $a\in M$ greater
than every standard numeral.  Interpret each fixed scalar term in $M$
through its primitive recursive defining function.  In particular, an
iterator with a nonstandard counter is interpreted by its internal
primitive recursive function in $M$, not by an external loop of
nonstandard length.

Let $\overline{k}$ denote the numeral in $M$ for a standard natural number
$k$.  Define the external subset
\begin{equation}
 U=\{x\in M:\text{ for some }k\in\mathbb N\text{ in the metatheory, }
                      M\models x<J(\overline{k},a)\}.
 \label{eq:finite-menu-cut}
\end{equation}
The quantifier over $k$ in this definition is an \emph{external standard}
quantifier.  It is not a quantifier over all elements of $M$, and no
definability of $U$ within $M$ is asserted.

The following properties follow from PA theorems already established.
First, $U$ is an initial segment of $M$.  It contains all standard numerals,
and it contains $a$, because $a<H(a)=J(1,a)$.  Second, it is closed under
successor: if $x<J(\overline{k},a)$, then
\[
 Sx\le J(\overline{k},a)
      <J(\overline{k+1},a).
\]
Every nonzero member $y$ of $U$ has a predecessor $z$ in $M$; since $z<y$,
that predecessor belongs to $U$.  Consequently every member of $U$ is
either $0$ or the successor of another member of $U$.

Third, $U$ is closed under every operation in $F$.  Given finitely many
arguments $x_1,\ldots,x_{r_i}\in U$, choose one standard $k$ large enough
that all of them are less than $J(\overline{k},a)$.  Put
$n=J(\overline{k},a)$ in \eqref{eq:finite-menu-majorant}.  Then
\[
 f_i(x_1,\ldots,x_{r_i})<H(J(\overline{k},a))
                       =J(\overline{k+1},a).
\]
Nullary operations satisfy the same bound without a list of arguments.
It follows by induction on an ordinary finite target expression that
every member of $\mathcal T_F(V)$ has value in $U$ whenever the parameters
in $V$ have values in $U$.

Finally, $g(a)$ does not belong to $U$.  For every standard $k$,
$M\models\overline{k}<a$, and \eqref{eq:finite-menu-orbit} gives
\[
 J(\overline{k},a)<J(a,a)=g(a).
\]
Thus $g(a)$ is below none of the bounds in
\eqref{eq:finite-menu-cut}.  This proves that $U$ is a proper cut, although
it is closed under all the permitted specialization targets.

\subsection{The two-carrier counterreading and rule soundness}

Terms take values in the full carrier $M$.  Numerical eigenparameters and
the variables bound by $\mathsf Q$ and $\mathsf D$ range over $U$.  For
precision, define the value $\operatorname{val}_{\rho}(t)\in M$ and the
truth value of a formula under \emph{any} $M$-valued assignment $\rho$ to
its free variables.  The numerical term interpretation is the one just
specified.  Write $A_\rho$ for the truth of $A$ under $\rho$, and write
$B_\rho(c,b)$ for its truth after assigning $x=c$ and $y=b$; the one-variable
notation is analogous.  Formula truth is defined by
\begin{align*}
 (s=t)_\rho
 &\quad\Longleftrightarrow\quad
   \operatorname{val}_{\rho}(s)=\operatorname{val}_{\rho}(t),\\
 (\mathsf Q_yB(y))_\rho
 &\quad\Longleftrightarrow\quad
   \bigl(B_\rho(0)\Longrightarrow\forall b\in U\,B_\rho(b)\bigr),\\
 (\mathsf D_{x,y}B(x,y))_\rho
 &\quad\Longleftrightarrow\quad
   \exists c\in U\,
   \bigl(B_\rho(c,0)\Longrightarrow\forall b\in U\,B_\rho(c,b)\bigr).
\end{align*}
In these displays, substitution of a member of $M$ means modifying the
assignment, rather than inserting a new syntactic constant.  A sequent is
valid when it preserves truth for all assignments of its free numerical
parameters to $U$.

This is a \emph{two-carrier reading of the specified inference rules}.
It is deliberately not an ordinary first-order structure for the complete
term language on domain $U$: $U$ is not closed under every term operation,
as witnessed by $g(a)$.  Nor is it the usual first-order reading on domain
$M$, since the logical binders range over $U$.  Soundness must therefore
be proved rule by rule; it is not being inferred from ordinary
first-order soundness.

The extension of the definition to arbitrary $M$-valued assignments gives
the usual substitution lemma for numerical terms and for native formulas.
In particular, replacing a term by a term of equal $M$-value preserves the
truth of any native formula.  The logical binders still range over the
same fixed set $U$ after such a replacement.  The iterator's internal
accumulator, in contrast, ranges over $M$ through its numerical evaluation.

\paragraph{Lemma (soundness of $\mathcal N_F$).}
\emph{Every derivable sequent of $\mathcal N_F$ is valid in the reading
above.}

\begin{proof}
Induct on the derivation.  Hypothesis is immediate, and both computation
equations hold in $M$.  Equality transport follows from the substitution
lemma, regardless of whether the two equal term values lie in $U$.

For specialization, the induction hypotheses give the closure and its
zero instance.  The chosen target has value in $U$, by the closure property
proved above.  The required instance is therefore one of the instances
asserted by that closure.

For elimination of $\mathsf D$, its true major premise supplies some
$c\in U$ for which the opened $\mathsf Q$-formula is true.  Assign the fresh
parameter $b$ the value $c$ and use the induction hypothesis for the second
premise.  Freshness ensures that the truth values of the undischarged
context and conclusion have not changed.

For introduction by cases, the chosen witness $w$ has value $c\in U$,
because it is $0$ or an existing parameter.  Fix an assignment making
$\Gamma$ and $B(c,0)$ true.  If $b\in U$ is $0$, the desired instance is
the assumed base.  Otherwise $b=Sz$ for a $z\in U$.  Assign that value to
the fresh parameter of the premise.  The induction hypothesis yields
$B(c,Sz)=B(c,b)$.  Thus the conditional closure holds at the permitted
witness $c$.

It remains to check equation progress.  Fix a parameter assignment making
$\Gamma$ true, and suppose $E(0)$ is true.  The induction hypothesis for
the premise says externally that
\[
 \text{for every }z\in U,\qquad E(z)\text{ true }\Longrightarrow E(Sz)
 \text{ true}.
\]
For any fixed $b\in U$, every $z<b$ in $M$ belongs to $U$, because $U$ is
an initial segment.  As $E$ is a literal equation, its truth is that of a
fixed internal arithmetic formula of $M$, with its numerical parameters
fixed.  Hence
\[
 M\models E(0)\ \wedge\ \forall z<b\,(E(z)\longrightarrow E(Sz)).
\]
PA proves the bounded induction consequence
\begin{equation}
 \left(E(0)\wedge\forall z<b\,(E(z)\longrightarrow E(Sz))\right)
       \longrightarrow E(b).
 \label{eq:finite-menu-bounded-induction}
\end{equation}
Applying this theorem inside $M$ yields $E(b)$.  Since $b$ was arbitrary
in $U$, the required $\mathsf Q$-closure is true.

No induction involving a predicate for $U$ was used.  The undischarged
context may have external truth conditions involving $U$, but after fixing
an assignment satisfying that context, the matrix in
\eqref{eq:finite-menu-bounded-induction} is only the internal equation
$E$.  This is exactly where the literal-equation restriction on progress
is used.
\end{proof}

The same argument handles the witnessed presentation of equation progress:
first fix its permitted witness in $U$, then apply the last paragraph to
the resulting equation.  Imposing additional syntactic restrictions on
transport, on specialization, or on the choice among the permitted
witnesses cannot invalidate this soundness proof.

\subsection{A closed native theorem that is lost}

For numerical terms $p,q$ define the selector
\[
 C(y,p,q):=I(y,p;u.q),
\]
with $u$ fresh, and define the native naming formula
\[
 \operatorname{Name}(t)
 :=\mathsf D_{x,y}\bigl(C(y,t,x)=t\bigr),
\]
where $x,y$ do not occur free in $t$.

\paragraph{Lemma (the naming test).}
\emph{Under any assignment in the counterreading,
$\operatorname{Name}(t)$ is true if and only if
$\operatorname{val}_{\rho}(t)\in U$.}

\begin{proof}
Write $p=\operatorname{val}_{\rho}(t)$.  If the naming formula is true,
choose its witness $c\in U$.  Its zero instance is true because
$C(0,p,c)=p$.  Its instance at $1\in U$ must therefore be true, and
$C(1,p,c)=c$ gives $c=p$.  Hence $p\in U$.

Conversely, if $p\in U$, select $c=p$.  For every $y\in M$ one has
$C(y,p,p)=p$: the zero case is the zero computation, and for $y=Sz$ the
successor computation returns the constant body $p$.  PA supplies this
zero-or-successor distinction throughout $M$.  In particular all the
required instances for $y\in U$ are true, so the naming formula is true.
\end{proof}

Consider the closed $\mathsf Q$-free formula
\begin{equation}
 \Psi_g:=\mathsf D_{r,v}\operatorname{Name}(g(v)),
 \label{eq:finite-menu-lost-formula}
\end{equation}
where $r$ is unused.  Its base instance is true: $g(0)=J(0,0)=0$, and
$0\in U$.  Its instance at $v=a$ is false, because $g(a)\notin U$.
The witness $r$ cannot change either fact.  Thus $\Psi_g$ is false in the
counterreading and, by the soundness lemma, is not derivable in
$\mathcal N_F$.

The unrestricted native calculus has a direct proof of $\Psi_g$.
For any term $t$ and a parameter $w$ fresh for $t$, successor computation
proves
\[
 C(Sw,t,t)=t.
\]
Introduction by cases with the witness $t$ therefore proves
$\operatorname{Name}(t)$; the base assumption of that introduction is
unused.  In particular it proves $\operatorname{Name}(g(Sz))$ for a fresh
$z$.  Apply introduction by cases once more, now to the matrix
$\operatorname{Name}(g(v))$ with unused existential variable $r$ and
dummy witness $0$.  This gives \eqref{eq:finite-menu-lost-formula}.
Neither universal generalization nor an arithmetic completeness theorem is
needed for this native derivation.  The finite-menu obstruction follows.

\paragraph{A guarded public version.}
The failure can also be placed in the guarded public grammar.  For an
equation $E$ of the form $s=t$, put
\[
 \operatorname{Gate}(E)
  :=\mathsf D_{p,q}\bigl(C(q,0,s)=C(q,0,t)\bigr),
\]
where $p$ is unused and both bound variables are fresh.  This formula has
the same counterreading truth value as $E$.  Its zero row is $0=0$; its
row at $1\in U$ is $s=t$; and if $s=t$, all its rows are true.  The
unrestricted native calculus derives $\operatorname{Gate}(E)$ from $E$
by computation, transport, and introduction by cases.

Consequently the closed formula
\[
 \mathsf D_{r,v}\mathsf D_{x,y}
 \operatorname{Gate}\bigl(C(y,g(v),x)=g(v)\bigr)
\]
is also false in the counterreading and directly provable in the
unrestricted calculus: gate the computed successor equation in the naming
proof just given.  Its innermost matrix is literally of the guarded form
$I(q,0;u.h)=I(q,0;u.k)$.  This avoids attributing the obstruction solely to
an unguarded auxiliary goal.

\subsection{Consequences and limits}

There is an elementary counterreading for the further restriction in which
targets are only parameters and $1$, with other numerals disallowed as
targets.
Take $M=\mathbb N$ and $U=\{0,1\}$.  The same rule-soundness proof works:
cases on $U$ use only $0$ and $S0$, and equation progress over $U$ uses
only the step from $0$ to $1$.  Then
$\operatorname{Name}(S(S0))$ is already a false closed native theorem.
This finite reading does not allow successor-of-parameter targets, since
$U$ is not successor closed.  The nonstandard cut construction above is
what handles those targets, and much larger fixed menus, correctly.

The theorem applies to every finite menu of scalar numerical terms.  In
particular it applies to fixed addition, multiplication, exponentiation,
and any other represented primitive recursive operations.  Using the
established term-completeness of the scalar language, it can equivalently
be stated for an arbitrary finite menu of fixed primitive recursive
operations.  An arbitrary supplied iterator body is different: it gives
an unbounded family of target operations and is not covered by a fixed
finite menu.

This obstruction is compatible with separate completeness results for
restricted introduction witnesses and for restricted specialization
targets.  An unrestricted witness may leave $U$, making its introduction
rule unsound in the present reading.  An unrestricted specialization
target may also leave $U$, making its elimination rule unsound.  The
counterreading therefore applies precisely when the two sources of
freedom are restricted together.

Finally, the conclusion is a strict failure of the original native
consequence relation, including on closed guarded formulas.  Because the
reading is not an ordinary first-order structure on $U$, it must not be
repackaged as a theorem that every changed-domain interpretation is
impossible.  Such an interpretation would require its own translations,
soundness proofs, and definable composite isomorphisms.  None is ruled out
by the statement proved here.
